\documentclass[11pt,a4paper]{article}
\usepackage[T1]{fontenc}
\usepackage[utf8]{inputenc}
\usepackage[margin=2.5cm]{geometry}
\usepackage{enumitem}
\usepackage{parskip}
\usepackage{hyperref}
\usepackage{xcolor}

\title{Release Information for SR Automation}
\author{Otto Mierka}
\date{\today}

\begin{document}
\maketitle

This note summarizes the recent SSE execution updates that are relevant for embedding the new mechanisms into a
cluster-side automation sequence. The \textcolor{red}{red-highlighted} text addresses changes in the current
Gendie(/SSE) execution tracks and, where configuration entries are shown in red, those entries are to be provided
in the conventional simulation configuration file \texttt{setup.e3d}.

\section*{Background}

Previously, the production-oriented launcher logic was centered around \texttt{e3d\_start.py}. That path was mainly
preconfigured for momentum-equation CFD runs. For XSE/screw cases it could also handle temperature runs, but this
was not yet arranged as a clean production automation path.

This legacy path remains available. The original \texttt{e3d\_start.py} keeps its previous functionality, and the
previously used parameter files are still present so the established legacy SSE/gendie track can continue to be
used where backward compatibility is needed. In other words, the original track functions the same way as before.

The new setup introduces YAML-driven staged execution plans. This makes the solver sequence explicit while keeping
the runtime call itself a one-liner. In the FeatFloWer installation, the exported launcher and YAML helper files
are available in \texttt{tools/e3d\_scripts/}.

At the moment, the YAML-based update is mainly foreseen for manual executions. This already provides a practical
benefit for cases such as non-isothermal screw simulations.

Operationally, \texttt{e3d\_start.py} can be replaced one-to-one by \texttt{e3d\_start\_yaml.py} when the
corresponding execution track is supplied. For example, \texttt{e3d\_xse.yaml} selects the isothermal screw/XSE
case.

At present, no non-isothermal execution runs are requestable from SR, and no corresponding feature request has been
placed so far.

\section*{Main Improvement}

The user-facing execution is now reduced to a single command of the form:

\begin{flushleft}
\textcolor{red}{\texttt{python3 e3d\_start\_yaml.py -f <project-folder> -n <ranks> -y <yaml-plan>}}
\end{flushleft}

If the cluster workflow prefers \texttt{srun}, the same driver also supports the script option
\texttt{-u} / \texttt{--use-srun}.

Important clarification for cluster integration:

\begin{itemize}[leftmargin=*]
  \item In \texttt{e3d\_start.py} and \texttt{e3d\_start\_yaml.py}, \texttt{-u} and
  \texttt{--use-srun} are two spellings for the same script flag.
  \item This flag only switches the launcher from \texttt{mpirun ...} to \texttt{srun ...}.
  \item It is not the same as Slurm's own \texttt{srun -u} option.
  \item The driver does not append \texttt{-u} to the generated \texttt{srun}
  command; it only prefixes the solver call with \texttt{srun}.
  \item When \texttt{--use-srun} is active, the script still keeps its internal
  \texttt{-n <ranks>} value for logging and partitioning, but it does not translate that value into
  \texttt{srun -n <ranks>}. The actual task count must therefore come from the Slurm allocation or
  outer launch context.
\end{itemize}

\begin{verbatim}
python3 e3d_start_yaml.py -f <project-folder> -n 16 -y e3d_xse.yaml -u
\end{verbatim}

This means ``use the script's srun launch path''. It does not mean that the script
will generate \texttt{srun -u ...}, and it does not guarantee that \texttt{16} is forwarded as an
\texttt{srun -n 16} argument.

The important operational simplification is that the old manual hide/seek around the active
\texttt{q2p1\_param*.dat} files now happens behind the curtains inside the YAML driver. The YAML stages point to the
backing parameter files, and the driver copies them into the active \texttt{\_data/} locations required by the
solvers. This is easier to automate than the older production variant where the surrounding script logic had to
manage the parameter-file switching explicitly.

\section*{DIE Variants to Expose in Automation}

All four DIE variants are now available as predefined YAML plans.

\subsection*{1. Pure die run}

\begin{itemize}[leftmargin=*]
  \item YAML plan: \texttt{e3d\_die.yaml}
  \item Command pattern:
\end{itemize}

\begin{verbatim}
python3 e3d_start_yaml.py -f <project-folder> -n <ranks> -y e3d_die.yaml
\end{verbatim}

This is the isothermal single-material DIE path.

\subsection*{2. Die with multimat}

\begin{itemize}[leftmargin=*]
  \item YAML plan: \texttt{e3d\_die\_multi.yaml}
  \item Command pattern:
\end{itemize}

\begin{verbatim}
python3 e3d_start_yaml.py -f <project-folder> -n <ranks> -y e3d_die_multi.yaml
\end{verbatim}

This is the isothermal DIE path with alpha/material-distribution transport. It also resolves the previously used
multimaterial track with the long-criticized parameter-file ``switching game'' (\texttt{Huetchenspiel}), because the
required parameter activation is now handled by the YAML driver instead of by fragile external script logic.
Compared with the legacy track, the material-distribution step is no longer used only once to influence the CFD.
In the current YAML setup it is applied two times, once in \texttt{init} and once in \texttt{main}. This gives the
material distribution a stronger and more sustained coupling into the flow solution and allows runs to reach a
broader material-filled portion of the flow domain.

\subsection*{3. Die with temperature}

\begin{itemize}[leftmargin=*]
  \item YAML plan: \texttt{e3d\_die\_temp.yaml}
  \item Command pattern:
\end{itemize}

\begin{verbatim}
python3 e3d_start_yaml.py -f <project-folder> -n <ranks> -y e3d_die_temp.yaml
\end{verbatim}

This is the non-isothermal single-material DIE path.

\subsection*{4. Die with temperature and multimat}

\begin{itemize}[leftmargin=*]
  \item YAML plan: \texttt{e3d\_die\_multi\_temp.yaml}
  \item Command pattern:
\end{itemize}

\begin{verbatim}
python3 e3d_start_yaml.py -f <project-folder> -n <ranks> -y e3d_die_multi_temp.yaml
\end{verbatim}

This is the non-isothermal multi-material DIE path. An important practical point is that the heat solver is no
longer limited to the pure CFD path. With the staged YAML setup it can be combined both with the CFD solver alone
and with the multi-material module in one normalized execution pipeline. As in the isothermal multi-material YAML
track, the material-distribution step is applied two times, once in \texttt{init} and once in \texttt{main},
rather than influencing the CFD only once as in the older legacy handling.

\section*{Why This Matters for Cluster Automation}

The relevant benefit is not only that these combinations exist, but that they are now normalized behind one
entrypoint.

\begin{itemize}[leftmargin=*]
  \item The selection of the physical run variant is reduced to choosing the YAML file.
  \item The concrete solver ordering is encoded in the YAML stage definition instead of in external ad hoc shell logic.
  \item The parameter-file switching is internalized in the driver.
  \item The same launcher syntax can be used across the different DIE variants.
  \item The driver performs a YAML preflight check and stops early if referenced parameter files are missing.
  \item Stage outputs are archived into legacy-compatible \texttt{\_prot<stage-index>} folders, while the protocol
        files from the \texttt{final} stage are collected in \texttt{\_prot\_final}. For report generation and other
        post-run analysis, the relevant protocol set should be taken from \texttt{\_prot\_final} so that the final
        stage is used consistently.
\end{itemize}

In short, the new path is more suitable for robust production sequencing than the older \texttt{e3d\_start.py}
workflow. At the same time, it should be seen as an additive upgrade rather than a removal of the old path, since
the legacy launcher and legacy parameterization remain available.

\section*{YAML Stage Progress in \texttt{e3d.log}}

The YAML driver now also exposes the staged execution state through \texttt{e3d.log}. This is intended for
cluster-side monitoring tools that already inspect the \texttt{[SimulationStatus]} section and should not have to
parse terminal output to understand where a staged YAML run currently is.

The status block now uses YAML-specific keys:

\begin{verbatim}
YamlStageProgress=<current-stage>/<total-stages>
YamlInStageProgress=<current-solver-in-stage>/<solvers-in-stage>
CurrentYamlSolver=<MomentumEquation|HeatEquation|MaterialDistribution>
CurrentYamlSolverProgress=<solver-local-current>/<solver-local-maximum>
CurrentYamlStatus=<status-text>
\end{verbatim}

When the YAML driver is active, transient status keys are intentionally named with \texttt{Yaml} in the keyword.
For example, \texttt{CurrentStatus} is emitted as \texttt{CurrentYamlStatus}, and the final timestamp is emitted as
\texttt{YamlFinishingTime}. Legacy inner-iteration keys such as \texttt{CurrentInnerIteration} and
\texttt{MaxInnerIteration} are not repeated in the YAML status block because the same information is represented by
\texttt{CurrentYamlSolverProgress}.

The total stage count is computed from the active YAML plan:

\begin{itemize}[leftmargin=*]
  \item \texttt{init} contributes one stage if it is present and has steps.
  \item \texttt{main} contributes its configured \texttt{loop} count.
  \item \texttt{final} contributes one stage if it is present and has steps.
\end{itemize}

For example, a plan with:

\begin{verbatim}
stages:
  init:
    steps:
      - solver: MomentumEquation
      - solver: MaterialDistribution
      - solver: HeatEquation

  main:
    loop: 2
    steps:
      - solver: MomentumEquation
      - solver: MaterialDistribution
      - solver: HeatEquation

  final:
    steps:
      - solver: MomentumEquation
\end{verbatim}

has four YAML stages in total:

\begin{verbatim}
1/4 = init
2/4 = main loop 1
3/4 = main loop 2
4/4 = final
\end{verbatim}

Inside each stage, \texttt{YamlInStageProgress} tracks normalized stage work, not just the raw solver index. This
matters for XSE/screw runs because one \texttt{MomentumEquation} step can expand to multiple angular positions. If
a momentum solve needs \texttt{L} angular positions, it contributes \texttt{L} in-stage progress units.
\texttt{HeatEquation} and \texttt{MaterialDistribution} each contribute one in-stage progress unit.

For an \texttt{init} or \texttt{main} stage with
\texttt{MomentumEquation + MaterialDistribution + HeatEquation}, the denominator is therefore \texttt{L + 2}. If an
XSE run needs \texttt{18} angular positions, the momentum solve advances through \texttt{1/20}, \texttt{2/20}, ...
\texttt{18/20}, the material-distribution step is \texttt{19/20}, and the heat step is \texttt{20/20}. For a DIE
single-angle run, \texttt{L = 1}, so the same stage naturally becomes \texttt{1/3}, \texttt{2/3}, and
\texttt{3/3}.

The solver-local progress is sampled from the solver protocol file where possible. For momentum and heat solves,
the driver watches \texttt{\_data/prot.txt} and parses the existing \texttt{itns: current/max} lines. For the
material-distribution solve, the driver watches the Q1 scalar output and maps the reported
\texttt{VolumeFractions[\%]} filling degree to an integer progress value against \texttt{100}.

Representative \texttt{e3d.log} status snapshots for the example pipeline are shown below. The surrounding fixed
header lines such as \texttt{PathToE3DFile}, \texttt{NumOfCpus}, and \texttt{YamlStartingTime} are omitted here for
brevity. In YAML mode, the legacy header entries \texttt{NumOfAnglePositions} and \texttt{Periodicity} are not
emitted; angular progress is represented through \texttt{YamlInStageProgress},
\texttt{CurrentYamlAngleIteration}, and \texttt{MaxYamlAngleIteration}.

During the \texttt{init} momentum solve:

\begin{verbatim}
YamlStageProgress=1/4
YamlInStageProgress=4/20
CurrentYamlSolver=MomentumEquation
CurrentYamlSolverProgress=4/12
CurrentYamlAngleIteration=4
MaxYamlAngleIteration=18
CurrentYamlStatus=running Momentum Solver
\end{verbatim}

During the first \texttt{main} material-distribution solve:

\begin{verbatim}
YamlStageProgress=2/4
YamlInStageProgress=19/20
CurrentYamlSolver=MaterialDistribution
CurrentYamlSolverProgress=67/100
CurrentYamlMaterialFillingDegree=67
MaxYamlMaterialFillingDegree=100
CurrentYamlStatus=running Material Distribution Solver
\end{verbatim}

During the second \texttt{main} heat solve:

\begin{verbatim}
YamlStageProgress=3/4
YamlInStageProgress=20/20
CurrentYamlSolver=HeatEquation
CurrentYamlSolverProgress=6/10
CurrentYamlStatus=running Heat Solver
\end{verbatim}

At the end of the \texttt{final} momentum stage:

\begin{verbatim}
YamlStageProgress=4/4
YamlInStageProgress=18/18
CurrentYamlSolver=MomentumEquation
CurrentYamlSolverProgress=12/12
CurrentYamlStatus=finished
YamlFinishingTime=...
\end{verbatim}

This keeps YAML-mode monitoring unambiguous: automation can restrict itself to keys containing \texttt{Yaml} and
still distinguish between \texttt{init}, repeated \texttt{main} stages, and \texttt{final}.

\section*{Constant-Mesh Heat Optimization}

There is one further runtime-relevant improvement that should be visible to automation users.

If the setup provides:

\begin{flushleft}
\textcolor{red}{\texttt{[E3DSimulationSettings]}}\\
\textcolor{red}{\texttt{ConstantMesh = yes}}
\end{flushleft}

These \textcolor{red}{additional entries are to be provided in the conventional simulation
configuration file \texttt{setup.e3d}}.

Then the code detects that the mesh is constant and treats the heat transport accordingly. In that case the matrix
set for the heat solve is built once initially and then reused, instead of being rebuilt repeatedly during the run.

Operationally, this matters because constant-mesh SSE-TEMP style runs avoid unnecessary matrix regeneration while
still using the same staged execution framework.

\section*{Relation to Screw/XSE Cases}

The same YAML concept also exists for screw/XSE cases:

\begin{itemize}[leftmargin=*]
  \item \texttt{e3d\_xse.yaml}
  \item \texttt{e3d\_xse\_temp.yaml}
\end{itemize}

So the DIE rollout is not a one-off mechanism. It fits into a broader staged-launch concept for SSE-related
executions.

\section*{Optional Extension: ROUND Geometry Keyword}

There is a second, independent update that may be relevant for automation as an optional feature flag.

If the setup contains:

\begin{flushleft}
\textcolor{red}{\texttt{[E3DGeometryData/Machine]}}\\
\textcolor{red}{\texttt{GeometryType = ROUND}}
\end{flushleft}

This \textcolor{red}{additional entry is likewise to be provided in the conventional simulation
configuration file \texttt{setup.e3d}}.

Then the new-generation mesher reacts to this keyword. That mesh generation step happens in the separate meshing
module, not inside FFer itself.

The intended workflow is:

\begin{enumerate}[leftmargin=*]
  \item The geometry setup advertises \texttt{GeometryType = ROUND}.
  \item The new mesher creates the corresponding round mesh.
  \item FFer then prolongates that mesh to higher levels using cylindrical parametrization for the higher-level
        nodes/vertices.
\end{enumerate}

This means the keyword is not just descriptive. It activates a geometry-dependent path spanning meshing and later
in-solver prolongation.

\section*{Boundary Between Mesher and FFer}

\begin{itemize}[leftmargin=*]
  \item Mesh creation for the round setup is handled by the separate, recently rolled-out meshing module.
  \item The higher-level coordinate prolongation is handled on the FFer side.
  \item The FFer prolongation uses round/cylindrical information when building the higher mesh levels.
\end{itemize}

So cluster automation may need to treat \texttt{GeometryType = ROUND} as a keyword-level switch that affects both
the meshing stage and the subsequent simulation startup path.

\section*{Summary Message}

The SSE execution path now supports four predefined YAML-driven DIE pipelines
(\texttt{die}, \texttt{die\_multi}, \texttt{die\_temp}, \texttt{die\_multi\_temp}) through a single launcher,
\texttt{e3d\_start\_yaml.py}. This reduces automation to selecting the appropriate YAML plan while the driver
handles internal parameter-file activation and stage sequencing, including the former multimat parameter-file
switching track. The heat solver can now be embedded both in pure CFD runs and in combined CFD plus
multi-material runs through the same mechanism. If \texttt{ConstantMesh = yes} is provided, the heat path also
avoids repeated matrix rebuilds by reusing the initially assembled transport matrices. In addition,
\texttt{GeometryType = ROUND} can optionally activate the newer round-mesh plus cylindrical-prolongation workflow,
with mesh creation in the separate mesher and higher-level coordinate prolongation in FFer. The original
\texttt{e3d\_start.py} path and the older parameter files remain available as a legacy-compatible fallback for the
traditional SSE/gendie workflow.

\end{document}
